\documentclass[10pt,letterpaper]{article}
\usepackage{cornell-notes}

\title{Clause 26.05.07.03: Ranges To Adaptors}
\author{}
\date{}
\setDocTitle{Clause 26.05.07.03: Ranges To Adaptors}
\setDocSubtitle{C++ 2024}
\setDocOwner{C++ 2024 Cornell Notes}
\setDocDate{\today}
\setCornellCollection{C++ 2024 Cornell Notes}
\setCornellUnitType{Clause}
\setCornellUnitNumber{26.05.07.03}
\setCornellUnitTitle{Ranges To Adaptors}
\hypersetup{pdftitle={Clause 26.05.07.03: Ranges To Adaptors},pdfsubject={Cornell notes},pdfauthor={}}

\begin{document}
\maketitle
\begin{CornellOverviewBox}{Overview}
\textbf{Classification:} Standard library - Normative\\[2pt]
\textbf{Learning objective:} Understand the rules, terminology, and semantic consequences associated with ranges::to adaptors.
\end{CornellOverviewBox}\begin{CornellSummaryBox}{Summary}
{\sffamily\bfseries\color{Accent} Summary}\par\vspace{3pt}Section 26.5.7.3 focuses on ranges::to adaptors. Apply the specified interface together with its mandates, effects, result conditions, exception behavior, and complexity constraints. When applying it, preserve the distinction between mandatory requirements, recommendations, permissions, and behavior deliberately left undefined, unspecified, or implementation-defined.
\end{CornellSummaryBox}

\end{document}
